\section*{Περίληψη στα Ελληνικά}
Η εργασία συγκρίνει παραδοσιακούς πράκτορες παιχνιδιών βασισμένους σε κανόνες
(Finite State Machine σε συνδυασμό με τον αλγόριθμο A*) με πράκτορες Ενισχυτικής
Μάθησης (Reinforcement Learning), εκπαιδευμένους με τον αλγόριθμο PPO, σε ένα
συμμετρικό, ανταγωνιστικό 2D περιβάλλον που αναπτύχθηκε στο Unity/ML-Agents.
Στο περιβάλλον αυτό, δύο πράκτορες μονομαχούν σε μια αρένα με δυναμικά,
τυχαία παραγόμενα εμπόδια (procedural generation), υπό ένα σύστημα πόντων ζωής,
βλημάτων και ελέγχου θερμοκρασίας πυρός.

Εξετάστηκαν πέντε πειραματικά σενάρια Ενισχυτικής Μάθησης με διάφορες μετρικές:
\begin{itemize}
	\item PPO εναντίον της απλής (baseline) εκδοχής του rule-based πράκτορα, σε 5 και 10 εκατομμύρια βήματα,
	\item PPO εναντίον της σύνθετης (advanced) εκδοχής του rule-based πράκτορα,
	\item PPO με εκπαίδευση self-play,
	\item PPO με fine-tuning του self-play πράκτορα εναντίον της απλής εκδοχής του rule-based πράκτορα.
\end{itemize}

Τα αποτελέσματα υποδεικνύουν ότι η Ενισχυτική Μάθηση αποτελεί μία πρακτική και ανταγωνιστική λύση για την ανάπτυξη πρακτόρων σε περιβάλλοντα παιχνιδιών με ελάχιστη ανθρώπινη προσπάθεια.

\vspace{5cm}
\noindent {\fontsize{14}{14}\selectfont \textbf{Λέξεις κλειδιά:}} ενισχυτική μάθηση, πράκτορες παιχνιδιών, αυτό-εκπαίδευση, πράκτορες βασισμένοι σε κανόνες

\pagebreak
\section*{Abstract ή Περίληψη στα Αγγλικά}

This thesis compares traditional rule-based game agents, built on a Finite
State Machine combined with the A* pathfinding algorithm, against
Reinforcement Learning agents trained with the Proximal Policy Optimization
(PPO) algorithm, within a symmetric, adversarial 2D combat environment
developed in Unity using ML-Agents. The environment features procedurally
generated, destructible obstacles, a health and resource-based shooting system,
and two agents competing under identical rules.

Several training approaches were evaluated, including PPO trained against
static rule-based opponents of varying difficulty, PPO trained through
self-play, and a subsequent fine-tuning phase applied to the self-play agent.
The results show that RL agents are capable of consistently outperforming
rule-based opponents while spontaneously developing behaviors typically
associated with hand-crafted agents, such as positioning, cover usage, and
projectile avoidance, without any explicit behavioral scripting and without
resorting to a time-consuming reward shaping process. Self-play training
further encouraged the emergence of active, generalizable targeting skills,
while a short fine-tuning phase against a static opponent substantially
improved both performance and generalization at minimal additional
computational cost.

Overall, the findings suggest that Reinforcement Learning constitutes a
practical and viable alternative to conventional rule-based approaches for
game agent development, capable of producing competent, adaptable agents
with comparatively little manual effort.

\vspace{8cm}
\noindent {\fontsize{14}{14}\selectfont \textbf{Keywords:}} reinforcement learning, PPO, game agents, self-play, rule-based agents, unity